\documentclass[11pt]{article}

\usepackage[T1]{fontenc}
\usepackage{lmodern}
\usepackage{amsmath}
\usepackage{amssymb}
\usepackage[margin=1in]{geometry}
\usepackage{booktabs}
\usepackage{hyperref}

\newcommand{\key}[1]{\texttt{#1}}

\title{Biased sampling in the Monte Carlo module:\\ algorithms, parameters and when to use them}
\author{Athena++ Monte Carlo radiation transport}
\date{}

\begin{document}
\sloppy
\maketitle

\section{The estimator and what biasing may change}

Every escaping sample contributes its weight $w$ to the spectrum. The run estimates a
band flux as $F = \sum_{\rm esc} w$ and its statistical error as
$\sigma_F = \sqrt{\sum_{\rm esc} w^2}$, both sums over the samples landing in the band.
The sampling is unbiased when the expectation of $\sum w$ over all escaping samples equals
the true escaping luminosity, whatever the individual weights are.

Every feature described here changes \emph{where samples are drawn and how often they
interact}, and compensates by changing their weights so that the expectation is untouched.
None of them changes the answer; all of them change the error, the cost, or both. The
right way to compare two settings is the figure of merit of the band you care about,
\begin{equation}
  \mathrm{FoM} = \frac{1}{(\sigma_F/F)^2\, t_{\rm cpu}},
\end{equation}
which is what a setting buys per CPU second. The global effective sample count printed by
the run, $(\sum w)^2/\sum w^2$ over all escaping samples, is the same statistic for the
bolometric luminosity and is dominated by the heaviest samples wherever they land in the
spectrum; it is a check for pathology, not a measure of the hard band.

\section{Sample allocation: \key{weights}}

\key{<montecarlo> weights} selects how emission cells are drawn.
\begin{description}
\item[\key{emission}] cells are drawn uniformly and the cell's emission goes into the
  weight. Cheap to set up and the weights follow the emissivity, which varies by orders
  of magnitude across a disk.
\item[\key{equal}] cells are drawn in proportion to their emission; every sample starts
  with the same weight $E/N$, the total emission over the number of samples. This is the
  natural analog scheme and the reference for everything below.
\item[\key{biased}] cells are drawn in proportion to emission times a per-cell importance
  $I$, and the weight is divided by $I$: $w = (S/N)/I$ with $S = \sum_c E_c I_c$. Cells
  the importance rates highly get more samples, each lighter. Unbiased for any positive
  $I$.
\end{description}
The biased scheme is worth its setup when most of the emission happens where it cannot
escape. In the X-ray binary disk over $99.9\%$ of the emission sits in cells of Thomson
depth above 30 per cell, and the equal-weight run spends most of its samples on photons
that are absorbed after hundreds of scatterings. With the escape-probability importance
the samples are moved to the skin, and the soft band ($<0.3$~keV) gains a factor of three
in error for twice the CPU.

\section{The importance: escape probabilities}

The importance the \key{mc\_readhdf\_gr} generator uses is an estimate of the probability
that a photon emitted in the cell escapes without being absorbed. It is built from the
effective extinction of a scattering, absorbing medium,
$\kappa_{\rm eff} = \sqrt{3\alpha_a(\alpha_a+\alpha_s)}$, the inverse of the
thermalization length, so that the cell's escape probability at energy group $l$ is
\begin{equation}
  p_l = \max\!\left(e^{-\tau_{l,\min}},\ \key{pesc\_min}\right),
\end{equation}
where $\tau_{l,\min}$ is the smallest column of $\kappa_{\rm eff}$ from the cell to the
domain boundary along the six axis directions. The cell's own contribution is not its
full depth but $-\ln[(1-e^{-\tau_c})/\tau_c]$, the escape probability of a sample placed
uniformly in it, without which a thick emitting skin cell would be rated unable to emit
at all.

The cell importance is the spectrum-weighted mean of $p_l$ over the groups,
\begin{equation}
  I = \frac{\sum_l s_l\, p_l}{\sum_l s_l},
\end{equation}
with $s_l = e^{-x}\Delta\ln E$ for free-free emission ($x = E/kT$) and the table's photon
shares for table emission.

\paragraph{Where the columns come from.} Without a file the columns are integrated within
each block only, which overstates escape in the face layer of every block, one cell in
eight on a 16-cell block. \key{vis/python/montecarlo/escape\_table.py} integrates them
through the whole mesh and writes the per-group probabilities on the snapshot's block
structure; \key{<problem> escape\_file} reads that file. Use the file for any production
run; the block-local estimate is for tests.

\paragraph{Parameters.}
\begin{center}
\begin{tabular}{@{}llp{3.1in}@{}}
\toprule
key & default & meaning \\
\midrule
\key{<problem> nescape} & 32 & energy groups of the free-free table (table emission uses its own grid) \\
\key{<problem> pesc\_min} & $10^{-10}$ & floor on $p_l$; bounds the weight a deep sample can carry at $(S/N)/\key{pesc\_min}$ \\
\key{<problem> escape\_file} & none & global escape table from \key{escape\_table.py} \\
\bottomrule
\end{tabular}
\end{center}
\key{pesc\_min} rarely needs changing. A deep sample with $I = \key{pesc\_min}$ carries
an enormous weight but almost never escapes; the composite mixing of the cell draw
(\key{bias\_mix}) and of the energy draw (\key{bias\_energy\_mix}) is the tool that
bounds it, not the floor.

\paragraph{What this importance is and is not.} It is the probability of escaping at all,
which is the right importance for the soft flux that leaves the skin directly. It is not
the probability of escaping \emph{in a given band}: in the X-ray binary the photons above
10~keV are emitted at $3$--$10\times10^6$~K a few scattering depths in, and scatter
400--600 times before leaving. Those cells have $I \sim 10^{-3}$, so their samples carry
$10^3$ times the skin weight, and the few that come out hard dominate the hard-band error.
The weight window (Section~\ref{sec:window}) is the correction for that; a
band-specific importance would need a tally from a previous run and is not implemented.

\section{Energy biasing in \key{PhotonEmitFreeFree}: \key{bias\_energy}}
\label{sec:energy}

The free-free sampler draws $\ln E$ uniformly over the cell's range ($x$ from
\key{tfloor} to \key{tceiling} in units of $kT$) and puts the spectrum into the weight,
$w \propto e^{-x}\,\Delta\ln E$. This deliberately oversamples the Wien tail relative to
its share of the emission, by $e^{x}$, so that a hard band has samples at all.

With an escape table and \key{bias\_energy = true} (the default in the biased scheme),
group $l$ is drawn with probability proportional to its overlap $d_l$ with the cell's
range times
\begin{equation}
  m_l = (1-\xi_E) + \xi_E\,\frac{p_l}{\langle p\rangle},
  \qquad \langle p\rangle = \frac{\sum_l d_l p_l}{\sum_l d_l},
\end{equation}
and the weight measure becomes $\Delta\ln E / m_l$ instead of the flat $\Delta\ln E$.
$\xi_E$ is \key{<montecarlo> bias\_energy\_mix}, which follows \key{bias\_mix} unless
set. Groups that can escape get more samples, each lighter; the expectation is
unchanged; and no sample's measure exceeds $1/(1-\xi_E)$ times the flat one, which is
the same bound the composite mixing of the next section puts on the cell weight. The
table sampler of \key{mc\_readhdf\_gr} does the same with the table's photon shares in
place of $d_l$. This helps where absorption is strongly energy dependent, as free-free
absorption ($\propto \nu^{-3}$) is: the soft groups of a skin cell are absorbed and the
hard ones escape, and the biased draw stops spending samples on the former.

\paragraph{Why the bound is needed.} At $\xi_E = 1$ the measure is
$\langle p\rangle\Delta\ln E/p_l$, and a group the table holds at the \key{pesc\_min}
floor is drawn $10^{10}$ times less often than one that escapes freely. The table rates
a sample by its birth energy in its birth cell and cannot see what scattering does to
it afterwards. On the X-ray binary that is exactly the class the hard bands are made of:
photons born at 0.1--0.25~keV in $10^6$~K gas that come out at 0.3--10~keV after
Compton upscattering, from cells where the table has the 32--100~eV band at the floor
in 13\% of cases and the 100--316~eV band there in 3\%. In expectation they are still
counted, but through samples so rare that a run of $10^7$ photons never draws one,
and when one does arrive it is a giant (the 2941-copy family of
Section~\ref{sec:window} was one). The biased run was 8--24\% low from 0.3 to 10~keV
against the emission run for this reason, with no other difference between the inputs
(Section~\ref{sec:xrb}). With $\xi_E = 0.9$ the suppression is at most tenfold and the
class is sampled. On an emission-only sphere (no absorption or scattering, so the
escaping spectrum is the emission spectrum) with a synthetic table whose soft bands sit
at the floor, $\xi_E = 0.9$ reproduces the equal-weight total and band shares to 0.3\%
over four outputs of $10^6$ samples, while $\xi_E = 1$ comes out 9\% low with 30--70
effective samples per output.

Set \key{bias\_energy = false} to keep the cell allocation but draw energies as the
other schemes do. That is the right choice when the escape probabilities are nearly flat
in energy (pure scattering, or grey absorption), where the biased draw buys nothing and
only adds weight spread, and it is the way to isolate the effect of the cell allocation
in a test.

\section{Composite mixing: \key{bias\_mix}}

Pure importance sampling gives a cell of importance $I$ the weight $(S/N)/I$, which for a
deep cell is unbounded above. Baes et al.\ (2016) mix the importance allocation with the
equal-weight one: the sampling density becomes $(1-\xi)\,E_c/E + \xi\,E_c I_c/S$, so the
weight of any cell is
\begin{equation}
  w_c = \frac{E/N}{(1-\xi) + \xi\, I_c E/S} \le \frac{1}{1-\xi}\,\frac{E}{N}.
\end{equation}
\key{<montecarlo> bias\_mix} is $\xi$; the default $0.99$ bounds every weight at 100
equal weights; $1$ is pure importance sampling. The code applies the mix once, at the
first initialization, by replacing the stored importance with $(1-\xi) + \xi I E/S$,
so everything downstream (window centres included) sees the mixed value; later outputs
of a \key{nout > 1} run reuse the mixed array and the reference weight recorded then.

\paragraph{What it does and does not fix.} The bound acts on samples from cells whose
$I E/S$ is below $1-\xi$: in the X-ray binary that is $I < 1.6\times10^{-7}$ at
$\xi = 0.99$, the deep disk. Those samples are 1\% of the run, carry 100 equal weights,
and the rare one that escapes lands in the 0.3--4~keV range, where seven such photons
carried 8\% of the escaped luminosity in one run. Lowering $\xi$ to $0.9$ cuts their
weight tenfold at the cost of 10\% of the samples drawn at equal weight, and an estimate
from the photon list puts the gain at a factor of two in effective samples in the
1--3~keV band and nothing elsewhere. The hard-band outliers come from cells with
$I \sim 10^{-3}$, for which $I E/S \gg 1-\xi$ at any reasonable $\xi$; the mix cannot
reach them.

\paragraph{Choosing $\xi$.} $0.99$ when the run is dominated by the soft band; $0.9$ when
the mid band matters; never below $0.5$, where half the samples go back to the analog
allocation and the soft-band gain is lost. The reference weight for \key{minweight} is
the importance-only average $S/N$, so it does not move with $\xi$.

\section{Minimum weight and roulette: \key{minweight}, \key{roulette}}

A sample whose weight falls below \key{minweight} times the reference weight ($E/N$ in
the equal scheme, $S/N$ in the biased one, the maximum cell emission in the emission
scheme) plays Russian roulette at its next interaction: it survives with probability
\key{roulette} (default $0.1$) carrying $1/\key{roulette}$ times its weight, or is
absorbed. This is unbiased and lets \key{minweight} be set high, $10^{-4}$ say, where
the old absorb-outright rule would have biased the result; the point is to stop
following samples whose weight has decayed under \key{abs\_method = weight} to where
they cannot matter. With the weight window's floor on this roulette is never reached,
since the window acts at a fraction of the birth weight rather than at $10^{-4}$ of it;
it remains the only roulette in the emission and equal schemes.

\section{Weight window: \key{wwin\_top}, \key{wwin\_bottom}, \key{wwin\_max\_split}}
\label{sec:window}

Importance sampling at emission gives the ideal weight only while a sample stays where it
was born. A sample that random-walks from a cell of importance $10^{-3}$ to one of
importance $0.1$ now carries a hundred times the weight of its neighbours; one walking
the other way costs as much to push as ever while its expected contribution has become
negligible. The window corrects both at every interaction:
\begin{itemize}
\item the sample's reference is the cell's window centre, $(S/N)$ over the mixed
  importance at the sample's current energy group, times the sample's own ratio of weight
  to centre at emission, stored in \key{Photon::wrefp};
\item above \key{wwin\_top} times the reference the sample is split into
  $\min(\lceil w/\mathrm{ref}\rceil, \key{wwin\_max\_split})$ equal copies at the same
  state, which the transport loop scatters independently;
\item below \key{wwin\_bottom} times the reference it survives at the reference with
  probability $w/\mathrm{ref}$ or is absorbed carrying no weight.
\end{itemize}
Both bounds default to 0, which disables that side; the window needs
\key{weights = biased} for its importance. The run report counts copies and roulettes.

\paragraph{Why relative to the birth ratio.} A window set per cell alone would compare
every sample with the cell's average weight. The energy sampler draws $\ln E$ flat and
carries $e^{-x}$ in the weight, so Wien-tail samples are born far below any per-cell
centre; such a window rouletted them wholesale and collapsed the 10--30~keV band of a
test sphere to a single survivor, while splitting the soft samples at their first
interaction. Referenced to each sample's birth ratio, the window acts only on the change
since emission, migration between cells and energy groups and absorption decay, and
leaves the emission's choices alone.

\paragraph{Where the reference is taken.} The centre is taken at the sample's
current energy group (\key{wwin\_energy}, default true): per cell, a hard photon
entering a cooler cell is split although its own escape prospects did not change, and
on a Compton sphere that made fifty times the copies. It is interpolated in $\ln p$
between cell centres to the sample's position (\key{wwin\_interp}, default true): a
cell-centred reference jumps by the neighbouring cells' importance ratio at every face,
15 in the X-ray binary disk, while a random walker moves a fiftieth of a cell per
scattering, so it is split and rouletted at every crossing, a fifth of all interactions.
Interpolation halved the events on the sphere.

\paragraph{Copies are correlated.} The copies of one birth photon share its history to
the split, so the photon-level effective count and the quoted errors of a window run
overstate the statistics badly: on one X-ray binary output a single birth photon's 2941
escaping copies carried 27\% of the luminosity and 64\% of the 1--3~keV band, 22,745
effective photons but 13 effective families. \key{family\_stats.py} groups the
escaping photons by birth photon through the emission temperature and emitted energy
columns and prints both counts, and \key{make\_spectrum.py --family-cols 0 1} writes
family-level errors, the root of the sum of squared family totals per bin, which is the
correct estimator with copies and reduces to the usual one without them; on the output
above it raised the 1--3~keV error from 4\% to 34\%. The scatter between outputs of a
\key{nout > 1} run is the error that needs neither.

\paragraph{The floor amplifies importance errors.} A sample the table rates unlikely to
escape is rouletted as it walks inward and its survivor carries $1/\key{wwin\_bottom}$
per round; if the table is wrong and such samples do escape, a family that survived $n$
rounds comes out with $4^n$ times the weight. Since the minimum-column escape
probability underestimates diffusive and oblique escape by orders of magnitude, this is
the origin of the giant families. On the cool Compton sphere the split alone reached
26\% of the equal-weight figure of merit, both sides 11--13\%, and the biased scheme
without a window 0.7\%. Start with \key{wwin\_top = 4}, \key{wwin\_bottom = 0}, and
turn the floor on only against a measured importance.

\paragraph{Choosing the bounds.} Copy counts approaching the number of interactions
mean the window is churning; widen it or raise \key{wwin\_max\_split} only if the split
is what is churning. The window is inert in a problem where samples do not migrate in
importance: a thin medium, or one where the importance is nearly uniform.

\section{Path stretching: \key{stretch}, \key{stretch\_taucell}, \key{stretch\_bound}}

Free paths are drawn with the extinction multiplied by $\alpha = \key{stretch}$, and the
weight carries the ratio of true to stretched survival, $e^{(\alpha-1)\tau}$ over the
depth flown since the last interaction, with a further factor $1/\alpha$ at the
interaction itself. A sample crossing a thin layer of depth $\tau$ then interacts with
probability $1 - e^{-\alpha\tau}$ instead of $1 - e^{-\tau}$, at proportionally lower
weight. The scheme is exact for any $\alpha$ that depends on position and on the
sample's history, which is what allows the two controls:
\begin{itemize}
\item \key{stretch\_taucell} (default unlimited) applies the multiplier only in cells
  whose depth across their smallest width is below the value, so a thick disk stays
  analog and its cost does not rise by $\alpha$;
\item \key{stretch\_bound} (default 10) caps the weight factor one flight may gain,
  after which the flight continues analog.
\end{itemize}
\key{Photon::strp} holds the flight's log weight factor and \key{Photon::taup} the depth
left in the flight, carried across blocks so that a flight whose interaction point lies
on a neighbouring block interacts there. Before \key{taup} existed such interactions were
redrawn on the new block and thereby skipped, about 1.6\% of them on a 16-cell block; with
stretching that was a bias of several percent, and in analog transport it delayed those
scatterings. The fix applies to every scheme.

\paragraph{When it helps.} When the observable comes from photons that cross a thin
scattering layer once or twice: a corona of $\tau \sim 0.3$ over a cool source, where
the analog run wastes most samples on photons that pass through unscattered. In a sphere
of $\tau = 0.9$ inside a $\tau = 0.3$ shell at $10^9$~K, \key{stretch = 2} confined to
the shell halved the Comptonized-tail errors for 8\% more CPU. The gain compounds over
the scatterings that make a tail photon, so larger values lose: 3 gained nothing more and
5 tripled the errors, and applying it to the sphere as well as the shell was worse still.

\paragraph{When it does not.} When the hard photons are made by long random walks in
thick hot gas rather than by single crossings of a thin layer. In the X-ray binary the
hard band is built by 400--600 scatterings in cells of Thomson depth 30--100, no thin
corona is crossed, and \key{stretch = 2} on the cells below depth 0.03 changed the
hard-band errors only by noise for 6\% more CPU. Leave it off unless the geometry has a
genuine thin layer and the target photons cross it a few times.

\section{Reading the diagnostics}

The run report prints, per output interval, the escaped weight and its effective sample
count, the counts of escaped, absorbed and removed samples, the mean scattering count,
and when active the window's copy and roulette counts and the stretching settings.
\begin{itemize}
\item The escaped weight must agree between settings within its error; the scheme is
  unbiased and any drift of several sigma is a bug.
\item The global effective count is set by the heaviest samples. In the X-ray binary it
  moved from 6372 to 1263 between two runs of the same settings because one deep-cell
  photon escaped in one and not the other. Compare per-band errors, from the spectrum
  file or from the photon list, and divide by CPU.
\item A photon list with the user columns (emission temperature, emitted energy,
  scattering count) tells where the heavy samples come from, which decides which tool
  applies: mixing for deep-cell outliers, the window for climbers, stretching for
  crossers.
\item The error $\sqrt{\sum w^2}$ is itself a sample estimate, and with heavy-tailed
  weights it is biased low: the samples that would dominate it are the ones the run
  did not draw. On a cool, dense sphere five seeds of the same biased-window run
  scattered by 1.5\% while each quoted 0.3\%, and the hundred heaviest escaping
  photons carried 20--50\% of $\sum w^2$; the equal-weight run quoted 0.24\% and
  scattered by 0.2\%. A quoted error is trustworthy only when no small group of
  photons dominates $\sum w^2$. When one does, run with \key{nout > 1} and take the
  scatter between the outputs' spectra as the error, which assumes nothing about the
  weights; \key{combine\_spectra.py -m time} averages them afterwards.
\item The quoted error can also come out too large. The emission scheme hands every
  block the same number of samples, so its sampling is stratified by block, and
  $\sum w^2$, which assumes independent draws, overstates the variance: on the X-ray
  binary the emission run's quoted errors were 1.4 to 3 times the scatter between its
  four outputs. The scatter is the measurement; the quoted error is a bound whose
  direction depends on the scheme.
\item Runs that share \key{iseed} draw the same cells and are not independent. Two
  emission runs with the same seed agreeing to 2\% says nothing about their errors,
  and a shape difference between schemes that recurs across such runs is not thereby
  confirmed. Change the seed when a repeat is meant as a check.
\end{itemize}

\section{What the X-ray binary runs showed}
\label{sec:xrb}

The schemes were compared on the high-resolution snapshot with four outputs of
$2.5\times10^6$ samples each, so that every error is the scatter between outputs. The
biased run used the global escape table, \key{bias\_mix = 0.9}, the split-only window
(\key{wwin\_top = 4}, \key{wwin\_bottom = 0}) with the interpolated reference, and no
stretching; the emission run was the standard input. CPU: biased $3.85\times10^6$~s,
emission $2.46\times10^5$~s, a ratio of 15.7.

\begin{center}
\begin{tabular}{@{}lcccc@{}}
\toprule
band [keV] & emission & biased & flux ratio & FoM biased/emission \\
\midrule
0.01--0.03 & 4.8\% & 4.1\% & 1.07 & 0.09 \\
0.03--0.1  & 2.3\% & 1.6\% & 1.13 & 0.14 \\
0.1--0.3   & 2.4\% & 0.6\% & 1.02 & 1.06 \\
0.3--1     & 2.1\% & 0.7\% & 0.93 & 0.69 \\
1--3       & 2.8\% & 1.6\% & 0.94 & 0.21 \\
3--10      & 6.8\% & 7.0\% & 0.82 & 0.06 \\
10--30     & 5.9\% & 6.9\% & 0.98 & 0.05 \\
\bottomrule
\end{tabular}
\end{center}

The biased run's errors are honest: scatter, family-level and photon-level estimates
agree, the family-level effective counts are 4000--6000 per output overall and
700--1500 in the 1--3~keV band, and no family carries more than 1\% of an output. Its
absolute errors are smaller in the soft and middle bands, by four times at 0.1--0.3~keV,
but that is what sixteen times the CPU would buy the emission scheme anyway, and in the
hard bands, where the hard photons are made by hundreds of scatterings in deep hot gas,
the split copies' walks buy nothing. Per unit CPU the biased scheme breaks even only at
0.1--0.3~keV. For this problem the emission scheme is the production tool.

The two schemes also disagreed on the spectral shape by more than either's scatter,
the biased run 6\% higher at 30--100~eV and 9, 8 and 24\% lower at 0.3--1, 1--3 and
3--10~keV. A second emission run with a different seed reproduced the first in every
band, and an equal-weight run could not arbitrate (0.06\% of its samples escape, 60
effective). Binning both runs' lists by birth temperature and emitted energy put the
whole deficit in photons born at 0.1--0.25~keV in $6\times10^5$--$10^6$~K gas that
escape two to ten times harder, which the energy draw at $\xi_E = 1$ with the
\key{pesc\_min} floor effectively never samples
(Section~\ref{sec:energy}); the emission run is right there, and
\key{bias\_energy\_mix} is the fix. The emission run's weights are the heavy-tailed ones
(130--600 effective photons per band against thousands of families for the biased run)
but that skews it low, not high, so it cannot be the source of the difference. The
biased run above was made before the mixed energy draw existed. A rerun with
\key{bias\_energy\_mix = 0.9} and $1.25\times10^6$ samples per output brought the total
onto the emission value and the 0.3--1 and 1--3~keV bands to within $1\sigma$ of it
after two outputs, at three times the old run's scatterings per output, since the
recovered photons are the ones that scatter hundreds of times; 3--10~keV stayed 20\%
low at $2\sigma$, where the emission run itself rests on about 60 effective photons per
output.

What the biased scheme would need here is an importance for the hard band, which the
escape probability is not. The cheap experiment not yet run is a floor at
\key{wwin\_bottom = 0.1} with the interpolated reference, which should fire only on
genuine inward migration and could cut the CPU several-fold; even so the scheme would
stay behind the emission run outside 0.1--1~keV. The machinery stays useful where it
fits: the family-level errors and the output scatter for any run, the carried flight
budget for all transport, and the window and stretching for problems with a genuine
thin scattering layer or an importance that varies gently from cell to cell.

\section{Settings that were tested on the X-ray binary}

These are the settings of the biased run above, kept as the reference configuration
for the biased scheme; see the previous section before choosing it over
\key{weights = emission} for this problem.

\begin{center}
\begin{tabular}{@{}llp{2.6in}@{}}
\toprule
key & value & why \\
\midrule
\key{weights} & \key{biased} & emission is deep; samples belong at the skin \\
\key{<problem> escape\_file} & from \key{escape\_table.py} & block-local columns overstate escape at block faces \\
\key{bias\_energy} & \key{true} & free-free absorption is strongly energy dependent \\
\key{bias\_energy\_mix} & 0.9 & bounds the energy weight factor at 10; at 1 the run above lost the soft-born upscattered photons \\
\key{bias\_mix} & 0.9--0.99 & 0.9 helps the 1--3~keV band; nothing below 0.5 \\
\key{minweight}, \key{roulette} & $10^{-4}$, 0.1 & harmless with the window; the only floor without it \\
\key{wwin\_top}, \key{wwin\_bottom} & 4, 0 & splits climbers; the floor amplifies table errors \\
\key{wwin\_max\_split} & 8 & default \\
\key{stretch} & 1 & no thin corona to cross \\
\bottomrule
\end{tabular}
\end{center}

And for a cool source under a thin hot corona: \key{weights = equal} or \key{biased} as
the source's depth demands, \key{stretch = 2}, \key{stretch\_taucell} between the
corona's and the source's cell depths, window on if biased.

\end{document}
